\documentclass[12pt]{article}

\usepackage{sbc-template}
\usepackage[brazil]{babel}
\usepackage[utf8]{inputenc}
\usepackage[T1]{fontenc}
\usepackage{graphicx}
\usepackage{url}
\usepackage{tikz}
\usepackage{booktabs}
\usepackage{array}
\usepackage{amsmath}
\usetikzlibrary{arrows.meta, positioning, shapes.geometric}

\sloppy

\title{Roteamento de Dados na Internet: Modelagem em Grafos e Comparação entre OSPF e RIP}

\author{Gabriel Pereira da Conceição\inst{1}, Kailan Marlon Spier\inst{2}, Rodrigo Colling Klaus\inst{3}}

\address{Escola Politécnica -- Universidade do Vale do Rio dos Sinos (UNISINOS)\\
  São Leopoldo -- RS -- Brasil
  \email{gabrielpereiradaconceicao@gmail.com}
}

\begin{document}

\maketitle

\begin{abstract}
  This article examines how the Internet's routing problem can be modeled as a
  weighted directed graph, where routers are vertices, network links are edges,
  and link costs are weights. Two widely deployed routing protocols are analyzed
  and experimentally compared: OSPF (Open Shortest Path First), which applies
  Dijkstra's algorithm to select minimum-cost paths, and RIP (Routing Information
  Protocol), which uses a distributed Bellman-Ford approach based solely on hop
  count. A virtual laboratory built with Containerlab and FRRouting demonstrates
  that over the same topology, OSPF correctly selects the higher-quality path while
  RIP chooses the path with fewer hops, which in this case has lower link quality.
  Failover behavior under link failure is also validated.
\end{abstract}

\begin{resumo}
  Este artigo analisa como o problema de roteamento na Internet pode ser modelado
  como um grafo ponderado, onde roteadores são vértices, enlaces são arestas e os
  custos dos enlaces são os pesos. Dois protocolos de roteamento amplamente
  utilizados são comparados experimentalmente: o OSPF (\textit{Open Shortest Path
  First}), que aplica o algoritmo de Dijkstra para selecionar caminhos de menor
  custo, e o RIP (\textit{Routing Information Protocol}), que utiliza uma variante
  distribuída do algoritmo de Bellman-Ford baseada apenas na contagem de saltos.
  Um laboratório virtual construído com Containerlab e FRRouting demonstra que,
  sobre a mesma topologia, o OSPF seleciona corretamente o caminho de maior
  qualidade enquanto o RIP escolhe o caminho com menos saltos, que neste caso possui
  enlace de menor qualidade. O comportamento de \textit{failover} após falha de
  enlace também é validado.
\end{resumo}

% =============================================================================
\section{Introdução}
% =============================================================================

A Internet é a maior rede de computadores já construída, interligando bilhões de
dispositivos por meio de uma malha de roteadores distribuídos por todo o mundo.
Para que um pacote de dados chegue ao seu destino, cada roteador no caminho precisa
decidir, de forma autônoma e em tempo real, para qual vizinho encaminhar o tráfego.
Esse processo é chamado de \textbf{roteamento}, e é regido por \textbf{protocolos
de roteamento} que os roteadores utilizam para trocar informações sobre a topologia
da rede e calcular os melhores caminhos \cite{tanenbaum2011}.

Dois dos protocolos mais relevantes nesse contexto são o \textbf{OSPF}
(\textit{Open Shortest Path First}) e o \textbf{RIP} (\textit{Routing Information
Protocol}). O OSPF é um protocolo de estado de enlace que constrói um mapa completo
da rede e aplica o algoritmo de Dijkstra para encontrar o caminho de menor custo
entre dois pontos. O RIP é um protocolo de vetor de distância que utiliza uma
variante distribuída do algoritmo de Bellman-Ford, tomando decisões com base apenas
na contagem de saltos — isto é, no número de roteadores traversados até o destino
\cite{rfc2328, rfc2453}.

A relação entre roteamento e teoria dos grafos é direta: uma rede de computadores
pode ser representada como um grafo ponderado, onde roteadores são vértices, enlaces
são arestas e os custos de cada enlace são os pesos. Nessa representação, calcular o
melhor caminho de roteamento equivale a resolver o problema do caminho de custo
mínimo em um grafo — problema clássico estudado há décadas na ciência da
computação.

Este trabalho tem como objetivo modelar formalmente o problema de roteamento como
um grafo ponderado, comparar o comportamento do OSPF e do RIP sobre a mesma
topologia e validar experimentalmente as diferenças por meio de um laboratório
virtual. A relevância prática é imediata: a escolha equivocada de protocolo ou
configuração incorreta de métricas pode degradar significativamente o desempenho de
uma rede de produção.

O artigo está organizado da seguinte forma: a Seção~\ref{sec:fundamentacao}
apresenta os conceitos teóricos de grafos e algoritmos envolvidos; a
Seção~\ref{sec:modelagem} descreve a modelagem do problema e a topologia utilizada;
a Seção~\ref{sec:discussao} apresenta os experimentos e discute os resultados; e a
Seção~\ref{sec:conclusao} sintetiza as conclusões.

% =============================================================================
\section{Fundamentação Teórica}
\label{sec:fundamentacao}
% =============================================================================

\subsection{Teoria dos Grafos e Roteamento}

Um \textbf{grafo} $G = (V, E)$ é composto por um conjunto de vértices $V$ e um
conjunto de arestas $E \subseteq V \times V$. Quando cada aresta possui um valor
numérico associado, diz-se que o grafo é \textbf{ponderado}. No contexto de redes
de computadores, os vértices representam roteadores e as arestas representam os
enlaces físicos ou lógicos entre eles. O peso de cada aresta pode expressar
latência, largura de banda, custo administrativo ou qualquer métrica relevante para
a qualidade do enlace \cite{cormen2022}.

O problema central do roteamento é encontrar o \textbf{caminho de custo mínimo}
entre um vértice de origem e um vértice de destino. Esse problema é resolvido de
formas distintas pelo OSPF e pelo RIP.

\subsection{Algoritmo de Dijkstra e o OSPF}

O algoritmo de Dijkstra resolve o problema do caminho mínimo de fonte única em
grafos com pesos não negativos com complexidade $O((V + E) \log V)$ utilizando uma
fila de prioridade \cite{dijkstra1959}. O OSPF adota uma abordagem de
\textit{link-state}: cada roteador propaga informações sobre seus próprios enlaces
para todos os outros roteadores da rede (via protocolo de \textit{flooding}),
permitindo que cada um construa um mapa completo e idêntico da topologia. Sobre
esse mapa, cada roteador executa o algoritmo de Dijkstra de forma independente para
calcular as rotas ótimas para todos os destinos \cite{rfc2328}.

A métrica padrão do OSPF é o \textbf{custo do enlace}, calculado em função da
largura de banda da interface. Custos menores indicam enlaces de maior qualidade;
o protocolo sempre prefere o caminho cujo somatório de custos seja o menor.

\subsection{Algoritmo de Bellman-Ford e o RIP}

O algoritmo de Bellman-Ford calcula caminhos mínimos a partir de uma única fonte em
grafos que podem conter arestas com pesos negativos, com complexidade $O(V \cdot
E)$ \cite{bellman1958}. O RIP implementa uma versão distribuída desse algoritmo,
conhecida como \textit{distance-vector}: cada roteador mantém uma tabela com a
distância estimada até cada destino e a compartilha periodicamente com seus
vizinhos diretos. Os roteadores atualizam suas tabelas ao receber as tabelas dos
vizinhos, convergindo iterativamente para as rotas ótimas.

A métrica do RIP é exclusivamente a \textbf{contagem de saltos} (\textit{hop
count}) — o número de roteadores atravessados até o destino. Esse critério ignora
completamente a qualidade dos enlaces, o que pode levar à seleção de caminhos de
desempenho inferior em redes heterogêneas. Adicionalmente, o RIP impõe um limite
máximo de 15 saltos, tornando-o inadequado para redes de grande escala \cite{rfc2453}.

% =============================================================================
\section{Modelagem do Problema}
\label{sec:modelagem}
% =============================================================================

\subsection{Representação em Grafo}

A rede foi modelada como um grafo ponderado não direcionado $G = (V, E, w)$, onde:

\begin{itemize}
  \item \textbf{Vértices} $V$: cada dispositivo de rede (roteador, switch, cliente
        e servidor) corresponde a um vértice;
  \item \textbf{Arestas} $E$: cada enlace físico entre dois dispositivos
        corresponde a uma aresta bidirecional;
  \item \textbf{Função de peso} $w: E \to \mathbb{Z}^+$: o custo OSPF do enlace,
        que reflete a qualidade relativa da conexão.
\end{itemize}

\subsection{Topologia e Objetivos do Design}

A topologia foi projetada para evidenciar três conceitos fundamentais: (i) o OSPF
prefere o caminho de menor custo total, mesmo que ele possua mais saltos; (ii) o
RIP seleciona o caminho de menor contagem de saltos, mesmo que ele possua enlaces
de menor qualidade; e (iii) o OSPF realiza \textit{failover} automático após falha
de enlace.

A Figura~\ref{fig:topologia} ilustra a topologia utilizada.

\begin{figure}[h]
\centering
\begin{tikzpicture}[
  node distance=1.8cm and 2.2cm,
  router/.style={circle, draw, thick, minimum size=1cm, font=\bfseries\small},
  host/.style={rectangle, draw, thick, minimum size=0.9cm, font=\small},
  link/.style={thick},
  ospf/.style={thick, blue!70!black},
  rip/.style={thick, red!70!black, dashed}
]
  \node[host] (client) {Client};
  \node[host, right=1.5cm of client] (s1) {S1};
  \node[router, right=1.8cm of s1] (r0) {R0};
  \node[router, above right=1.4cm and 2cm of r0] (r1) {R1};
  \node[router, below right=1.4cm and 2cm of r0] (r2) {R2};
  \node[router, right=2cm of r1] (r3) {R3};
  \node[host, below right=1.4cm and 1cm of r3] (s2) {S2};
  \node[host, right=1.5cm of s2] (server) {Server};

  \draw[link] (client) -- (s1);
  \draw[link] (s1) -- (r0);
  \draw[ospf] (r0) -- node[above, font=\tiny, blue!70!black]{custo 10} (r1);
  \draw[rip]  (r0) -- node[below, font=\tiny, red!70!black]{custo 40} (r2);
  \draw[link] (r1) -- (r3);
  \draw[link] (r3) -- (s2);
  \draw[link] (r2) -- (s2);
  \draw[link] (s2) -- (server);
\end{tikzpicture}
\caption{Topologia da rede. Caminho preferido pelo OSPF (azul, custo total 30):
  Client → S1 → R0 → R1 → R3 → S2 → Server.
  Caminho preferido pelo RIP (vermelho tracejado, 2 saltos):
  Client → S1 → R0 → R2 → S2 → Server.}
\label{fig:topologia}
\end{figure}

\subsection{Nós e Endereçamento}

A Tabela~\ref{tab:nos} descreve os nós da rede e seus papéis.

\begin{table}[h]
\centering
\caption{Nós da topologia}
\label{tab:nos}
\begin{tabular}{@{}lll@{}}
\toprule
\textbf{Nó} & \textbf{Imagem} & \textbf{Papel} \\
\midrule
client  & alpine:latest        & Origem do tráfego \\
server  & alpine:latest        & Destino do tráfego \\
s1, s2  & alpine:latest        & Switches de acesso \\
r0      & frrouting/frr:latest & Roteador core — ponto de entrada \\
r1      & frrouting/frr:latest & Caminho preferido pelo OSPF \\
r2      & frrouting/frr:latest & Caminho preferido pelo RIP / backup OSPF \\
r3      & frrouting/frr:latest & Ponto de saída via R1 \\
\bottomrule
\end{tabular}
\end{table}

A Tabela~\ref{tab:custos} apresenta os custos OSPF configurados nos enlaces.

\begin{table}[h]
\centering
\caption{Custos OSPF dos enlaces e seus significados}
\label{tab:custos}
\begin{tabular}{@{}llrl@{}}
\toprule
\textbf{Interface} & \textbf{Nó} & \textbf{Custo} & \textbf{Significado} \\
\midrule
eth2 & R0 & 10 & Enlace de alta qualidade em direção ao R1 \\
eth3 & R0 & 40 & Enlace de menor qualidade em direção ao R2 \\
eth1 & R2 & 40 & Custo simétrico no lado do R2 \\
demais & — & 10 & Custo padrão \\
\bottomrule
\end{tabular}
\end{table}

Com essa configuração, o custo total do caminho via R1 é $10 + 10 + 10 = 30$ e o
custo total do caminho via R2 é $40 + 10 = 50$. O RIP ignora esses custos e enxerga
o caminho via R2 como melhor por ter apenas 2 saltos contra 3 do caminho via R1.

% =============================================================================
\section{Discussão Técnica}
\label{sec:discussao}
% =============================================================================

\subsection{Ambiente de Laboratório}

O laboratório foi implementado com as seguintes ferramentas:

\begin{itemize}
  \item \textbf{Ubuntu 24.04 (WSL2)}: sistema operacional host;
  \item \textbf{Docker Engine}: \textit{runtime} de containers;
  \item \textbf{Containerlab 0.75.0}: orquestração da topologia de rede virtual,
        criando interfaces \textit{veth} entre containers para simular enlaces físicos;
  \item \textbf{FRRouting (FRR) 9.0}: suíte de protocolos de roteamento que
        implementa OSPF e RIP nos containers dos roteadores.
\end{itemize}

Os dois laboratórios — OSPF e RIP — foram executados em paralelo sobre a mesma
topologia lógica, com os mesmos endereços IP e a mesma estrutura de enlaces.

\subsection{Resultado 1 — Adjacência OSPF}

Após a inicialização do laboratório, todos os roteadores OSPF formaram adjacências
no estado \textit{Full}, indicando que o banco de dados de topologia (LSDB) estava
sincronizado e o algoritmo de Dijkstra havia sido executado.

\begin{verbatim}
Neighbor ID  Pri  State        Up Time  Address    Interface
1.1.1.1        1  Full/Backup  1m21s    10.0.1.2   eth2
2.2.2.2        1  Full/Backup  5m41s    10.0.2.2   eth3
\end{verbatim}

\subsection{Resultado 2 — Seleção de Rota: OSPF vs RIP}

A Tabela~\ref{tab:comparacao} compara a decisão de roteamento de cada protocolo
para o destino \texttt{10.0.5.0/30} (rede do servidor), extraída das tabelas de
roteamento de R0 em cada laboratório.

\begin{table}[h]
\centering
\caption{Comparação das decisões de roteamento em operação normal}
\label{tab:comparacao}
\begin{tabular}{@{}lll@{}}
\toprule
\textbf{Critério} & \textbf{OSPF} & \textbf{RIP} \\
\midrule
Algoritmo subjacente    & Dijkstra              & Bellman-Ford distribuído \\
Métrica utilizada       & Custo do enlace       & Contagem de saltos \\
Rota escolhida          & R0→R1→R3→S2           & R0→R2→S2 \\
Saltos                  & 3                     & 2 \\
Custo / métrica         & 30                    & 2 (saltos) \\
Caminho de maior qualidade selecionado & Sim   & Não \\
Limite de escala        & Sem limite prático    & Máximo 15 saltos \\
\bottomrule
\end{tabular}
\end{table}

O OSPF selecionou o caminho com três saltos e custo total 30, preterindo o caminho
com dois saltos mas custo 50. O RIP fez a escolha inversa: o caminho de dois saltos
foi preferido por ter métrica 2, enquanto o de três saltos tinha métrica 3 — sem
qualquer consideração sobre a qualidade dos enlaces.

Esse resultado demonstra empiricamente que o critério de seleção é determinante: em
redes heterogêneas, onde diferentes enlaces possuem capacidades distintas, o uso
exclusivo da contagem de saltos como métrica leva a decisões subótimas.

\subsection{Resultado 3 — Failover Automático no OSPF}

Para validar a capacidade de reconvergência, a interface \texttt{eth1} do roteador
R1 foi desativada manualmente, simulando uma falha de enlace:

\begin{verbatim}
docker exec clab-ospf-lab-r1 ip link set eth1 down
\end{verbatim}

Após o \textit{dead timer} do OSPF expirar e os LSAs de atualização serem
propagados, R1 foi removido da tabela de vizinhos e o protocolo recalculou as rotas.
O tráfego foi automaticamente redirecionado via R2 (custo 50), sem nenhuma
intervenção manual.

A Tabela~\ref{tab:failover} compara a tabela de roteamento antes e depois da falha.

\begin{table}[h]
\centering
\caption{Tabela de roteamento do R0 antes e após falha do enlace R0–R1}
\label{tab:failover}
\begin{tabular}{@{}lll@{}}
\toprule
\textbf{Destino} & \textbf{Operação normal} & \textbf{Após falha do R1} \\
\midrule
10.0.5.0/30 (servidor) & custo 30, via R1→R3→S2 & custo 50, via R2→S2 \\
Enlace R0\(\longleftrightarrow \)R1           & presente               & ausente (\textit{link down}) \\
\bottomrule
\end{tabular}
\end{table}

Ao restaurar a interface (\texttt{ip link set eth1 up}), o OSPF reconvergiu
novamente e reverteu ao caminho via R1 com custo 30, demonstrando que o protocolo
mantém o caminho ótimo enquanto a topologia estiver disponível.

Uma observação relevante emerge do failover: após a falha, o caminho ativo passou a
ter \textit{menos} saltos e \textit{maior} custo — reforçando que a métrica do OSPF
é sempre o custo, independentemente da contagem de saltos.

\subsection{Limitações e Desafios}

O RIP apresenta uma limitação estrutural relevante: o risco de
\textit{count-to-infinity} em topologias com rotas alternativas, onde atualizações
incorretas podem circular indefinidamente entre vizinhos até atingir o limite de 15
saltos \cite{rfc2453}. Esse problema não existe no OSPF, pois cada roteador possui
um mapa completo e consistente da topologia antes de calcular as rotas.

A ausência de suporte a múltiplas métricas no RIP é sua principal limitação
estrutural: o protocolo não possui mecanismo para considerar largura de banda,
latência ou confiabilidade dos enlaces na tomada de decisão de roteamento.

% =============================================================================
\section{Conclusão}
\label{sec:conclusao}
% =============================================================================

Este trabalho demonstrou que o problema de roteamento na Internet pode ser
naturalmente modelado como um grafo ponderado, onde algoritmos clássicos de caminho
mínimo têm aplicação direta. O OSPF, ao aplicar o algoritmo de Dijkstra sobre um
mapa completo da rede, seleciona consistentemente o caminho de menor custo total,
adaptando-se automaticamente a falhas de enlace. O RIP, por sua vez, utiliza uma
abordagem mais simples baseada em Bellman-Ford distribuído com contagem de saltos, o
que pode resultar em decisões subótimas em redes com enlaces heterogêneos.

Os experimentos realizados confirmaram quantitativamente as diferenças: sobre a
mesma topologia com custos deliberadamente assimétricos, o OSPF escolheu o caminho
com três saltos e custo 30, enquanto o RIP escolheu o caminho com dois saltos e
custo efetivo 50. O failover automático do OSPF também foi validado, com
redirecionamento do tráfego sem intervenção manual após a falha de enlace.

A relação entre a disciplina de Algoritmos e Estruturas de Dados e a infraestrutura
real da Internet é concreta: cada roteador em produção executa, a cada mudança na
topologia, algoritmos equivalentes aos estudados em sala de aula. Compreender essa
conexão é fundamental para projetar, operar e depurar redes de computadores.

% =============================================================================
\bibliographystyle{sbc}
\begin{thebibliography}{9}

\bibitem{tanenbaum2011}
Tanenbaum, A. S. e Wetherall, D. (2011).
\textit{Redes de Computadores}.
5ª edição. Pearson Prentice Hall.

\bibitem{cormen2022}
Cormen, T. H., Leiserson, C. E., Rivest, R. L. e Stein, C. (2022).
\textit{Introduction to Algorithms}.
4ª edição. MIT Press.

\bibitem{dijkstra1959}
Dijkstra, E. W. (1959).
A note on two problems in connexion with graphs.
\textit{Numerische Mathematik}, 1(1):269–271.

\bibitem{bellman1958}
Bellman, R. (1958).
On a routing problem.
\textit{Quarterly of Applied Mathematics}, 16(1):87–90.

\bibitem{rfc2328}
Moy, J. (1998).
\textit{OSPF Version 2}.
RFC 2328. Internet Engineering Task Force (IETF).
Disponível em: \url{https://www.rfc-editor.org/rfc/rfc2328}

\bibitem{rfc2453}
Malkin, G. (1998).
\textit{RIP Version 2}.
RFC 2453. Internet Engineering Task Force (IETF).
Disponível em: \url{https://www.rfc-editor.org/rfc/rfc2453}

\end{thebibliography}

\end{document}
